\SourceLeaf{159}schung im Gerichtsarchiv ist ermittelt, dass die betreffenden jüngsten Kinder das Alter von 20 Jahren alle wirklich erreicht haben. Es kann indess sein, dass auf diese Weise immer noch etwas zu wenig gerechnet ist, da möglicherweise bei dem Tode eines oder des andern verwitweten Professors oder einer Professorwitwe noch Kinder zwischen 12 und 20 Jahren vorhanden gewesen sein können, die mithin nach jetziger Einrichtung pensionsberechtigt und also in meine Rechnung mit aufzunehmen gewesen sein würden: wenigstens haben die Gerichtsakten nicht in allen Fällen das Gegentheil zur Gewissheit gebracht. Indessen würde doch jedenfalls keine erhebliche Vergrösserung der Totalsumme dadurch hervorgebracht werden können.

3) Unter jenen 108 Theilnehmern befinden sich 7, deren Witwen noch jetzt am Leben sind. Für diese habe ich zu den Pensionszahlungen, welche sie bisher genossen haben, auch noch den schon oben nach der Wahrscheinlichkeitsrechnung evaluirten künftigen Betrag beigefügt, nachdem derselbe von der Epoche 1845 October 1 auf das Moment des Beitritts der resp. Ehemänner zurückdiscontirt war. Ohne dieses Verfahren hätte ich nicht bloss diese Interessenten ausschliessen, sondern consequenterweise anstatt der Theilnehmer 1—108 mich auf die 1—60 beschränken müssen, was ich jedoch hier nicht weiter entwickeln kann.

Zur Ersparung des Raumes setze ich die Resultate der 108 Rechnungen nicht einzeln hieher, sondern nur den Totalbetrag von allen. Es findet sich
\begin{center}
\footnotesize\setlength{\tabcolsep}{3pt}
\begin{tabular}{l r|r}
&\multicolumn{2}{c}{nach dem Zinsfuss}\\
&\(3\tfrac{1}{2}\) Proc.&4 Proc.\\\cline{2-3}
Summe der discontirten Beiträge . . . . .&13909,00 Thl.&13115,94 Thl.\\
—\quad —\quad —\quad Witwenpensionen&67915,80 ,, &59019,67 ,, \\
—\quad —\quad —\quad Waisenpensionen&10945,04 ,, &9878,34 ,, \\
Summe aller discontirten Pensionen . . .&78860,84 ,, &68898,01 ,,
\end{tabular}
\end{center}
\noindent\hspace*{\parindent}Es folgt hieraus zuvörderst, dass der (reducirte) Betrag der Waisenpensionen sehr nahe dem sechsten Theile des (reducirten) Betrages der Witwenpensionen gleich zu setzen ist: in der That ist jener bei dem Zinsfuss \(3\tfrac{1}{2}\) Proc. etwas kleiner, bei dem Zinsfuss 4 Proc. hingegen etwas grösser als dieses \(\tfrac{1}{6}\). Es liesse sich leicht nachweisen, dass diese kleine Verschiedenheit des Resultats nach Maassgabe des Zinsfusses vollkommen in der Natur der Sache begründet ist; da jedoch der Umfang der Erfahrungen zu klein ist, als dass dem Resultate überhaupt eine minutiöse Schärfe beigelegt werden könnte, so habe ich in der obigen Rechnung S. [152] mich schlechthin an das einfache Verhältniss \(1:6\) gehalten.

Man kann ferner sagen, dass die Gesammtrechte aller jener 108 Interessenten, nach jetziger Höhe der Pensionen und Beiträge taxirt und auf die Zeiten ihrer respectiven Eintritte reducirt, den Geldwerth hatten
\begin{center}
\begin{tabular}{l}
von 64951,84 Thalern bei dem Zinsfuss von \(3\tfrac{1}{2}\) Proc.\\
\phantom{von }55782,07 Thalern bei dem Zinsfuss von 4 Proc.
\end{tabular}
\end{center}
\noindent oder, so weit eine so ausgedehnte Erfahrung maassgebend sein kann, dass nach dem Durchschnittswerth, der Eintritt in die Kasse für jeden einzelnen den Geldwerth von 601,40 Thalern, oder von 516,50 Thalern hat, jenachdem man \(3\tfrac{1}{2}\) oder 4 Proc. Zinsen rechnet.

Mancher, der mit Gegenständen dieser Art wenig bekannt ist, wird sich vielleicht über die Kleinheit dieses Resultats wundern. Allein das Auffallende verschwindet grösstentheils, wenn man erwägt, dass dabei die Reduction auf den Zeitpunkt des Eintritts jedes einzelnen zum Grunde liegt, und dass jede Geldsumme, durch Discontirung für eine beträchtliche Anzahl von Jahren, enorm vermindert erscheint. Ich will diess durch das Beispiel von \textsc{Heyne} erläutern, welcher 48 Jahr Beitrag geleistet, und dessen Witwe 21 Jahr 11 Monat Pension genossen hat. Nach meinen Rechnungsgrundlagen stellt sich also die Summe der Beiträge zu 480 Rthl. und die Summe der Pensionsbezüge zu 5479 Rthl. 4 Ggr. heraus, mithin der
